\scanpage{46}
$A$ variété abélienne en car $p$. On a, pour la partie de $\pi_1(A)$ première à $p$
\[
\pi_1^{\neq p}(A)\simeq\varprojlim_{\substack{n\\\text{\unsure{arith}}\\(n,p)=1}}\deleted{\unsure{n}}A^{(n)}
\]
où $\deleted{\unsure{n}}A^{(n)}$ est le noyau de $A\xrightarrow{n}A$, et où $A^{(mn)}\mathbin{-}A^{(n)}$ est la multiplication par $m$.
\[
\boxed{\pi_1^\ell(A)\simeq\varprojlim_{\deleted{\unsure{\ell}}k} A^{(\ell^k)}\simeq T_\ell(A)=\operatorname{Hom}(\mathbb Q_\ell/\mathbb Z_\ell,A)}
\]
\margin{module de Tate}
$\ell$ étant un nb premier $\neq p$. Ce résultat provient du fait qu'un \uline{revêtement non ramifié d'une variété abélienne est une variété abélienne}

D'après Kummer, on a d'autre part
\[
\deleted{\text{Donc}}\operatorname{Hom}(\pi_1^\ell(A),U_\ell)\simeq
\text{sous-gpe de }\unsure{\mathfrak P|\mathfrak D}(A)=H^1(A,\underline O_A^{\unsure{*}})
\]
des élts d'ordre une puissance de $\ell$.

Si on admet que \uline{$A$ n'a pas de torsion}, on obtient donc
\[
\operatorname{Hom}(\pi_1^\ell(A),U_\ell)\simeq A^{*(\ell^\infty)},\qquad\text{i.e.}
\]
\[\boxed{T_\ell(A)\times T_\ell(A^*)\longrightarrow T_\ell(U_\ell)}\]
accouplement à valeurs dans le $\mathbb Z_\ell$-module libre de rang 1 $T_\ell U_\ell)$, et accouplement qui est une dualité.
\[
T_\ell(A)/\ell^kT_\ell(A)\times T_\ell(A^*)/\ell^kT_\ell(\unsure{A'|A^*})\longrightarrow
T_\ell(U_\ell)/\ell^kT_\ell(U_\ell)
\]
\[A^{(\ell^k)}\times A^{*(\ell^k)}\longrightarrow U^{(\ell^k)}\]
\margin{45}
